% ===== §2 Doctrinal Survey =====
\section{The Ethics of the Water-Virtues: A Doctrinal Survey}
\label{p20:sec:review}

\noindent
The water-virtues --- non-contention (\emph{Daodejing}, ch.~8), benefiting all things, dwelling in the low place, formlessness, and non-fullness --- are surveyed here across the modern, largely secular ethical traditions. The aim is not enumeration but \emph{location}: by letting the doctrines argue against one another over each virtue, the paper's own position is forced into view. Each virtue proves to be, in different doctrines, both affirmed and opposed; the survey's product is a set of standing challenges (\xref{p20:sub:review-summary}), handed forward to the dialectical core (\xref{p20:sec:dialectic}) and, in one case, answered in place. Throughout, the reading is structural rather than cultivational, for the reasons set out at the close (\xref{p20:sub:review-cultivation}): the water-virtues are treated as properties of how a relation is arranged, not as feats of an exceptional character.


% ---- Group A: political economy and the scale of distribution ----
\subsection{Political economy and the scale of distribution}
\label{p20:sub:review-polecon}

We begin with the traditions that read a relation not as a meeting of characters but as a circulation of value, because they are the ones that can give the first water-virtue --- benefiting without contending --- a body that owes nothing to self-cultivation. To benefit the ten thousand things and yet not contend (\emph{Daodejing}, ch.~8) sounds, in the moralist's ear, like a counsel of generosity: give freely, claim nothing. Read through political economy it becomes something sharper and less personal --- a claim about the \emph{form} of a productive relation, true or false of a structure regardless of the magnanimity of anyone inside it.

\subsubsection{Political-economic ethics}
\label{p20:ssub:review-marx}

The decisive reframing is to hear ``benefiting yet not contending'' as the description of a \emph{non-exploitative} relation. On the Marxian analysis, exploitation has an exact mark that is independent of anyone's intentions: value created by one party is appropriated by another, so that the surplus does not return to those whose activity produced it \parencite{marx1976}. Against this mark, ``benefiting the ten thousand things and not contending'' names its negation --- a relation in which what is generated flows back to the whole that generated it rather than being drawn off and held. The point of substituting this reading for the generous-character reading is that it relocates the good from the giver's heart to the relation's geometry. A relation can be exploitative while every party to it is kind, and non-exploitative while no one is especially generous; what settles the matter is where the surplus goes, not how it is felt. This is the first and most important gain of a structural reading: it makes the water-virtue a property a relation can have or lack, not a height a soul must climb to.

Two further traditions sharpen the same point from different sides. Sen's capability approach lets us state what ``benefiting'' should mean without collapsing it into welfare: to benefit another is to widen the range of lives genuinely open to them --- their capability set --- and not merely to raise the level of some good they consume \parencite{sen1999}. This matters because it rescues ``benefiting the ten thousand things'' from the utilitarian reading that would equate it with maximised satisfaction; one may raise another's welfare while narrowing their freedom, and that is not the water-virtue but its counterfeit. To benefit, in the sense the figure of water intends, is to nourish the other's own power to grow, not to administer a quantity of good to a passive recipient. Eglash's account of generative justice then supplies the return condition in a form already native to this series: value should cycle back to those whose generativity produced it, and the injustice of extraction is precisely the interruption of that return \parencite{eglash2016,eglash1999}. On this reading ``not contending'' is not passivity but the refusal to interpose a claim that would divert the cycle --- the refusal, in the language this series has used throughout, to convert circulating surplus into held stock.

The other water-virtue that political economy illuminates is dwelling in the low place. Read structurally, the low place is not a posture of humility but a \emph{position in a division of value}: the systematically devalued site of a relation, the labour or the need that the prevailing accounting does not see. Feminist political economy has named this site with particular precision in the case of reproductive labour --- the sustaining, repairing, attending work whose very invisibility is the condition of its exploitation \parencite{federici2004}. To dwell where others disdain is then not to abase oneself but to direct one's attention and one's flow toward exactly the position the structure has discounted. We will give this its own development later (\xref{p20:sec:lowliness}); here it is enough to register that the low place, like non-contention, is a structural coordinate before it is ever a virtue of temperament. Its connection to the rest of the series is direct: the devalued position is where the generative balance of \textcite{paperxvii} fails, where static inequality has been mistaken for settled fact, and the water-virtue is the flow that moves to correct it.

\subsubsection{Game-theoretic and economic ethics}
\label{p20:ssub:review-game}

If political economy supplies the water-virtue's content, game theory supplies its first and hardest test, and it is essential not to soften it. Stated as a disposition to cooperate without insisting on one's own claim, non-contention is, in a single encounter with a self-interested counterpart, simply the losing strategy: the one who does not contend is the one who is taken. In the one-shot prisoner's dilemma, unconditional cooperation is dominated; a flow that benefits all and contends for nothing is, against an appropriator, exploited to exhaustion. This is not a marginal objection but the precise point at which the generous-character reading collapses, and a structural reading must either answer it or concede that the water-virtue is a luxury available only where no one means harm.

The answer game theory itself supplies is that the verdict reverses once the encounter is not single but repeated, and the parties carry their histories. In the iterated setting, strategies that open with cooperation and are not exploitable --- that return cooperation for cooperation but withdraw it from persistent defection --- outperform unconditional appropriation over time, and cooperation becomes not a sacrifice but an equilibrium sustained by the shadow of the future \parencite{axelrod1984}. What converts non-contention from folly to rationality is therefore a definite structure: repetition, the legibility of the other's past conduct, and the capacity to make a commitment the other can believe \parencite{schelling1960}. This is a result of the first importance for the present paper, because it locates exactly what non-contention \emph{requires} in order not to be naive --- and what it requires is not a better character but a particular shape of relation, one in which acts are remembered and commitments can be made credible. That shape is the one this series has elsewhere called the non-binding vow: the betrothal structure whose whole work is to make a commitment credible without an external sword \parencite{paperxiii}. We mark the dependence here and discharge it later (\xref{p20:sub:dialectic-survival}).

This yields the first of the standing challenges the survey will accumulate.

\begin{claim}[Challenge 1, the rationality of non-contention]
Non-contention is rational, rather than self-defeating, only within a structure of repetition, reputation, and credible commitment. Absent such a structure it is dominated, and the water-virtue reduces to a disposition that survives only where it is not tested. Any ethics of non-contention therefore owes an account of the structure that makes it more than naivety.
\end{claim}

\noindent
The challenge is not answered by insisting that non-contention is noble; it is answered, if at all, by exhibiting the structure. We let it stand, and turn to the traditions of institutions and political philosophy, where the same demand returns in its sharpest form --- the claim that absent an enforcing power, non-contention is not merely irrational but suicidal.


% ---- Group B: political philosophy and the scale of institutions ----
\subsection{Political philosophy and the scale of institutions}
\label{p20:sub:review-polphil}

Political economy located the water-virtues in the circulation of value; political philosophy asks the prior question of how a relation of many parties is ordered at all --- by what claims, what recognitions, what powers. Here the water-virtues meet both their sharpest allies and their hardest opponent, and a structural reading must take the opposition as seriously as the support.

\subsubsection{Political ethics: Rawls, Honneth, Pettit}
\label{p20:ssub:review-polethics}

It must be said at once that the first virtue, non-contention, stands in genuine tension with the dominant modern theory of justice, and the tension should not be smoothed over. For Rawls, justice is what free persons would agree to precisely by \emph{pressing} their claims --- advancing competing demands from behind a veil that strips them only of the knowledge of their own position, not of the disposition to claim \parencite{rawls1971}. Justice on this account is the fair adjudication of contention, not its absence; a party who declined to press any claim would not be exhibiting a virtue but forfeiting their standing in the very procedure that secures fairness. This is a real objection to any ethics that would make non-contention foundational, and we register it rather than dissolve it: there are domains --- the basic structure of a society, the distribution of primary goods --- where the pressing of claims is not a failure of the water-virtue but the medium of justice itself, and where a premature non-contention would simply cede the field to those who do contend. The water-virtue is an ethics of intimate relation, and its scope is not automatically the scope of political justice; one of the paper's later tasks (\xref{p20:sec:watershed}) is to ask how far the dyadic virtue can scale, and Rawls marks one limit it may not cross.

Within that limit, however, two further traditions convert what looked like mere yielding into something with structural teeth. Honneth's theory of recognition gives the low place a precise content: the systematically misrecognised position --- the party whose contribution, suffering, or standing the prevailing order fails to see \parencite{honneth1995}. To dwell where others disdain is then to direct recognition toward exactly the site the order has withheld it from, and this is not self-abasement but a redistribution of recognition --- an active correction of a moral injury, not a lowering of oneself beneath one's due. Pettit's republicanism gives the complementary reading of ``benefiting yet not ruling.'' Its central good is non-domination: freedom not as the mere absence of interference but as the absence of any \emph{capacity} for arbitrary interference, any standing power of one over another \parencite{pettit1997}. ``To benefit and not rule'' --- to generate without establishing \emph{imperium} --- is exactly the relational form of non-domination: a position that gives without thereby acquiring a standing power to dispose of the one it gives to. This is the most secular possible rendering of the water-virtue, and it shows that ``not ruling'' is not weakness but the deliberate refusal to convert benefaction into dominion.

\subsubsection{Contractarianism: the hardest opponent}
\label{p20:ssub:review-hobbes}

The sharpest challenge to the water-virtue comes not from the theorists of fair contention but from the contractarian tradition that begins with Hobbes, and it must be stated without mitigation, because everything the dialectical core later does depends on having let it stand at full force. In the state of nature --- the condition of parties with no common power above them --- to benefit all and contend for nothing is not merely the losing move it was for game theory; it is fatal. Where there is no assurance that others will forbear, the one who forbears unilaterally delivers themselves to those who do not, and the rational course is the anticipatory contention Hobbes calls the war of all against all \parencite{hobbes1651}. Non-contention here is not a higher wisdom but an early death, and the figure of water seems to confirm rather than refute the charge: water does flow downhill, does yield, does take the lower place --- and the lower place is precisely where, on the Hobbesian picture, one is overrun.

This pushes the rationality challenge of Group A to its extreme, and we issue it as a distinct, graver challenge, because the structure it demands is correspondingly stronger.

\begin{claim}[Challenge 1$'$, the survival of non-contention]
Absent a power that makes others' forbearance credible, non-contention is not merely irrational but self-destroying: the one who does not contend is not exploited at the margin but eliminated. An ethics of the water-virtue therefore owes more than a reason why non-contention can be rational; it owes an account of how the water-virtue survives the absence of an external sword --- how it can be safe to be like water without first erecting a Leviathan to protect those who are.
\end{claim}

\noindent
We do not answer this here, and we resist the two easy answers. We do not say that intimate relations are simply not the state of nature, exempt by their warmth from the Hobbesian logic --- for the betrayal of intimacy is among the surest routes to exactly the predation Hobbes describes, and an ethics that assumed good faith would be no ethics for the cases that test it. Nor do we erect the sword: the entire burden of this series' treatment of the non-binding vow \parencite{paperxiii} is that intimacy's commitments are precisely those that hold \emph{without} an external enforcer, so that purchasing the water-virtue's safety with a Leviathan would forfeit what makes it the water-virtue. The challenge is therefore real and unpaid, and we carry it to the dialectical core (\xref{p20:sub:dialectic-survival}), where the answer will turn out not to be the abolition of the danger but the shaping of an ineliminable asymmetry --- the recognition, drawn from the formal companion \parencite{wan2026braid}, that the water-virtue never required a world without power, only a power whose surplus is returned rather than retained.

\subsubsection{Commons governance: Ostrom}
\label{p20:ssub:review-ostrom}

Between the theorists of contention and the theorist of the sword sits a body of work that answers both at the scale where the water-virtue most naturally lives --- not the dyad and not the state, but the shared resource held in common. Ostrom's study of common-pool resources showed, against the orthodoxy that a shared resource must be either privatised or policed from above, that communities sustainably govern shared waters, fisheries, and pastures through institutions of their own making: graduated commitments, mutual monitoring, locally adapted rules, and conflict-resolution that no external authority imposes \parencite{ostrom1990}. The example is not incidental to our figure: among Ostrom's central cases are literal watercourses --- irrigation systems whose users, over centuries, evolved the arrangements by which water benefits all of them and is contended for by none to the point of ruin. This matters to the present paper in two ways. First, it shows that ``benefiting all and not contending'' can be the stable property of an \emph{institution}, not a disposition of individuals --- the structural reading vindicated at the scale of the many. Second, it gives the paper's eventual move from the dyad to the watershed (\xref{p20:sec:watershed}) a genuine theory of governance rather than a metaphor: when we ask how the resonance of two can scale to a relational commons, Ostrom's design principles are the form the answer must take, and her irrigation communities are the proof that the water-virtue, institutionally embodied, is neither naive nor utopian. We mark the connection here and develop it where the scaling question is properly joined.


% ---- Group C: relational ontology, the many, and the ecological ----
\subsection{Relational ontology, the many, and the ecological scale}
\label{p20:sub:review-relational}

The traditions gathered so far read the water-virtues at the scale of two parties or of an institution. A third group reads them at the scale where water is no longer a figure but the thing itself --- the scale of the living whole, the basin, the web of mutual dependence --- and it is here that ``benefiting all things'' sheds the last of its sentimental sound and becomes a claim about how a self stands to the world that sustains it.

\subsubsection{Environmental and ecological ethics}
\label{p20:ssub:review-eco}

Leopold's land ethic performs, for the relation between a person and the living community, exactly the enlargement the water-virtue asks. Its single precept --- that a thing is right when it tends to preserve the integrity, stability, and beauty of the biotic community, and wrong when it tends otherwise --- extends the boundary of moral concern from persons to the whole of which persons are a part, so that the community itself, not any member, becomes the unit benefited \parencite{leopold1949}. This is ``benefiting all things'' stated as a structural obligation rather than a tender feeling: the good is what returns to the integrity of the whole, the bad is what draws the whole down for the gain of a part, and no appeal to anyone's character is needed to tell them apart. Naess's deep ecology sharpens the ontological point in a way that touches this series at its root. Against a shallow ecology that protects nature for the sake of human use, deep ecology holds that the self is not a bounded ego set over against its environment but is constituted by its relations to the wider field of life --- a relational self whose flourishing is not separable from the flourishing of what it is related to \parencite{naess1973}. This is, in the vocabulary of the founding text of the present series, the generative relational self \parencite{wan2024grb}: the individual as a knot of relation rather than a premise standing prior to it. The convergence is not decorative. It shows that ``benefiting all things'' and ``not contending'' are not demands laid upon a self from outside, against its interest, but expressions of what such a self already is: a being whose own good is bound up with the good of the field it belongs to, so that to benefit that field is not self-sacrifice but self-continuation. Callicott's systematic defence of the land ethic against the charge of a tyrannical holism --- his insistence that obligations to the whole are nested with, not destructive of, obligations to individuals --- supplies the qualification the water-virtue will also need \parencite{callicott1989}: that benefiting the whole must not become a licence to dissolve the particular, a caution we will meet again, in intimate form, as the danger of self-effacement (\xref{p20:sub:dialectic-bound}).

\subsubsection{Confucian ethics, read politically rather than as self-cultivation}
\label{p20:ssub:review-confucian}

Confucianism is the tradition most easily misread, for our purposes, as a discipline of self-cultivation --- the perfection of the self through ritual and the overcoming of the self's partiality. Read that way it would belong with the contemplative traditions we set aside (\xref{p20:sub:review-cultivation}). We read it instead for two of its structural doctrines, which give the water-virtue something the Daoist source conspicuously lacks: a graded order. The first is the graded extension of \emph{ren}, benevolence: the Confucian, against the Mohist doctrine of impartial concern for all alike, holds that care properly radiates outward in degrees --- from kin, to the people, to the things of the world --- so that to benefit all is not to benefit all \emph{equally} but to benefit each according to the concrete relation one stands in to them \parencite{confucius1998,mencius2003}. This is a genuine correction to ``benefiting all things,'' which in its Daoist statement can drift toward an undifferentiated levelling that, applied to intimacy, would be absurd: one does not stand to a stranger as to a spouse, and a love that benefited all alike would benefit no one as a beloved. The graded reading gives the water-virtue a \emph{structured} scale, an intra-Eastern tension that keeps benefiting all things from collapsing into a sentiment that honours everyone and binds to no one. The second doctrine is \emph{quan}, the weighing of the exigent against the canonical --- the recognition that the standing rule must sometimes be set aside for the situation, not by abandoning principle but by a higher fidelity to it \parencite{tu1985}. This is the Confucian form of ``lawful yet without fixed form,'' and it anticipates exactly the Kantian imperfect duty we will use to answer the charge that formlessness is lawlessness (\xref{p20:ssub:review-kant}): a norm that genuinely binds while leaving the form of its fulfilment to judgment. Confucianism, read for its structure rather than its disciplines, thus contributes the gradation and the situational judgment that the bare figure of water does not itself supply.

\subsubsection{Ubuntu relational ethics}
\label{p20:ssub:review-ubuntu}

The southern African ethic of Ubuntu states the relational ontology with a directness the other traditions approach only by argument: a person is a person through other persons --- one's very humanity is constituted in and by relation, not possessed prior to it \parencite{metz2011,mbiti1970}. Two of the water-virtues receive, in this light, a reading that removes the last trace of self-abnegation from them. Dwelling in the low place is not a descent beneath one's worth but an acknowledgement of a standing fact, that one's personhood is owed to the relations one is held in; there is no height to descend from, since the self was never the self-standing thing the objection presupposes. And benefiting others is not a transfer from a full self to an empty one but the maintenance of the relational fabric in which both parties are constituted --- so that to benefit the other is, quite literally and not metaphorically, to sustain the conditions of one's own personhood. Ubuntu thus rejoins, from an independent tradition, the conclusion deep ecology reached: that the water-virtues are not sacrifices demanded against the self's interest but expressions of a self that is relational through and through. We note the convergence with the generative relational lineage of this series, to which Ubuntu and Eglash's generative justice both belong, and carry forward its lesson --- that the appearance of self-sacrifice in the water-virtues is an artefact of the entity-first picture of the self that a relational ontology has already abandoned.


% ---- Group D: subject, duty, and the meta-ethical layer ----
\subsection{Subject, duty, and the meta-ethical layer}
\label{p20:sub:review-subject}

The remaining traditions do not enlarge the scale of the water-virtues but probe their interior: whether a formless good can be a law at all, what it does to the subject who practises it, and whether ``the water-good'' is the sort of thing that can be true. Two of the survey's five challenges are joined here, and one of them is answered in place.

\subsubsection{Kantian ethics: the high ground}
\label{p20:ssub:review-kant}

Kant is the natural adversary of an ethics of water, and confronting him squarely yields the survey's one outright victory rather than a deferral. The objection writes itself: morality, for Kant, is law, and a law must have form --- it must be a maxim one can will as a universal law, the same for all rational agents in like case \parencite{kant1998}. ``Take the shape of the vessel,'' ``flow wherever the ground leads,'' ``have no fixed form'': these seem the very negation of a moral law, a counsel to do whatever the situation invites, which is precisely what the categorical imperative was built to rule out. If the water-virtue is formlessness, then by Kant's lights it is not a lower kind of morality but the absence of morality, and the charge --- Challenge~3, that formlessness is lawlessness --- is the most principled of all the objections the survey collects.

The answer is to deny the premise that the water-virtue is formless, using a distinction internal to Kant's own system. Kant divides duties into perfect and imperfect \parencite{kant1996}. A perfect duty fixes the act completely: it specifies exactly what must be done or omitted, and admits no latitude --- the duty not to make a lying promise leaves nothing to discretion. An imperfect duty binds no less stringently as to \emph{whether} one must act, but leaves open the \emph{form} of the action: the duty of beneficence commands that one make others' ends one's own and help where one can, but it does not, and cannot, prescribe whom to help, when, by what means, or how much. It is a duty with a determinate ground and an indeterminate form --- lawful, universalisable, genuinely binding, yet without a fixed shape. This is exactly the structure of the water-virtue. ``Benefit all things'' is not the lawless ``do whatever the moment suggests''; it is the imperfect duty of beneficence, which obligates absolutely as to its end --- one must benefit, one may not be indifferent --- while leaving the channel of benefit to be found in the particular case, as water finds its course without thereby being formless. The formlessness that looked like lawlessness is revealed as latitude, and latitude is not the absence of law but a recognised mode of it. We have thus answered Challenge~3 with Kant's own resources, and without recourse to moral particularism --- without, that is, abandoning universal principle for the doctrine that the situation alone dictates the right. The water-virtue keeps the principle and locates its formlessness exactly where Kant himself locates the open form of an imperfect duty. The same move secures the reading of ``benefit yet do not rule'' as a case of the formula of humanity: to benefit another without ruling them is to treat them as an end and never merely as a means, the benefaction that does not convert the beneficiary into an instrument of the benefactor's will.

\subsubsection{Psychoanalysis: the danger of taking the receiver's shape}
\label{p20:ssub:review-psych}

Psychoanalysis contributes both a vindication and the survey's gravest warning, and the warning is what makes the dialectical core necessary. On the side of vindication, the Lacanian account of the imaginary lets us read non-contention not as moral restraint but as a release from a specific trap: the rivalrous identification of the imaginary register, in which the ego is constituted through a mirroring competition with its counterpart, so that to contend is to remain captured in that specular rivalry \parencite{lacan2006}. Non-contention, so read, is not the suppression of a desire to win but the dissolution of the very stage on which winning and losing are defined --- a withdrawal from the imaginary's zero-sum logic rather than a sacrifice within it. But the same framework names, with a precision no other tradition matches, the danger that the water-virtue courts. ``Taking the shape of the vessel'' is, read through the clinic, perilously close to the subject's self-erasure: becoming the shape the desiring Other wants, dissolving one's own desire into the demand of the other, which Lacanian ethics marks as the one thing not to be done --- \emph{céder sur son désir}, to give way on one's desire. The water that takes any vessel's shape may be the water-virtue; it may equally be the subject who has surrendered its desire and become a mere function of another's, and nothing in the figure of water, by itself, tells the two apart. This is Challenge~2 in its sharpest form, and we hand it forward with that sharpness intact (\xref{p20:sub:dialectic-bound}): an ethics of formlessness that cannot distinguish the generous fluidity of the water-virtue from the pathological fluidity of the effaced subject has not yet earned the name of an ethics. Klein supplies, alongside the warning, a corrective to one misreading of ``benefiting'': her account of reparation shows that the impulse to make good the damage one has done need not take the form of self-punishing sacrifice, the draining of the self to atone, but can be the creative transformation of the destructive impulse into something that builds \parencite{klein1937}. Benefiting, on this reading, is sublimation rather than self-immolation --- which is precisely the distinction the dialectical core must secure in general.

\subsubsection{Moral realism and the truth of the water-good}
\label{p20:ssub:review-metaethics}

A final, meta-ethical question must be settled, because the formal translation of the later sections depends on its answer. When we say that a relation has the water-good, or lacks it, are we describing a property the relation possesses independently of our attitude toward it, or merely projecting an approval that has no further truth-condition? If the latter, then the holonomy criterion the formal companion supplies (\xref{p20:sec:formal}) would be, at best, a precise encoding of a preference, and its air of objectivity a confusion. We therefore declare for a position, while noting it is a declaration the survey motivates rather than proves: a \emph{thin, relational moral realism}. The water-good is a structural property of the flow of value, power, and desire in a relation --- the sign and geometry of its circulating surplus --- and this property obtains, or fails to obtain, as a matter of how the relation is structured, not as a matter of how anyone feels about it \parencite{railton1986}. The realism is ``thin'' because it makes no claim that the good is a non-natural property or a feature of the universe apart from relations; it is \emph{relational} because the property it ascribes is precisely a property of relations, not of isolated states or acts. This is the natural meta-ethics for the structural reading announced at the outset (\xref{p20:sub:intro-structural}): if the water-good is a matter of the relation's geometry rather than the agent's cultivation, then it is as objective, and as discoverable, as that geometry --- and a formal model of that geometry is not an encoding of preference but a description of where the good, so understood, lies. It is this thin relational realism that licenses the translation layer to read the theorems of the companion paper as claims about the good, and not merely as claims about a formalism.


% ---- Group E: the cultivational reading set aside ----
\subsection{The cultivational reading set aside, and a methodological specimen}
\label{p20:sub:review-cultivation}

It remains to say explicitly what the survey has been steering around. There is a large and venerable body of thought in which the water-virtues are read as disciplines of the self: the emptying of self in Christian kenosis, the overcoming of selfish partiality in the Confucian \emph{keji}, the dissolution of self-grasping in Buddhist practice, the released equanimity of Heidegger's \emph{Gelassenheit} \parencite{heidegger1962}, the unselfing toward a quasi-divine condition that runs through the contemplative traditions east and west \parencite{nishida1990}. On these readings ``not contending,'' ``dwelling below,'' and ``emptying'' are exercises by which an exceptional soul climbs, against the grain of its self-interest, toward goodness. We have not drawn on them as the paper's frame, and the omission is deliberate rather than dismissive. Three reasons govern it. First, a cultivational reading locates the good in the agent's achieved interior, and so cannot distinguish, at the level of the relation, a generative descent from a dominative or a self-erasing one --- the very distinction the paper exists to draw (\xref{p20:sec:dialectic}); two relations may issue from equally disciplined souls and differ entirely in their geometry. Second, an ethics resting on exceptional attainment cannot be an ethics of ordinary intimate life, which is lived by people who are not saints and whose relations must be assessable without first assessing the sanctity of the partners. Third, and most pointedly, the cultivational reading is the one that makes the water-virtue sound like a counsel of self-abnegation, and it is exactly that sound the structural reading was introduced to silence: once the self is understood relationally (\xref{p20:ssub:review-eco}, \xref{p20:ssub:review-ubuntu}), the appearance of self-sacrifice dissolves, and with it the principal reason to read the virtues as ascetic feats.

That said, the contemplative material affords one positive service, as a specimen of the operation the whole survey performs. Stoicism is the instructive case, because it looks cultivational and is, at its core, structural. Its central counsels --- \emph{amor fati}, the love of one's fate; the acceptance of what is not up to us; the conforming of the will to the rational order of the whole --- can be heard as a spiritual discipline of resignation \parencite{marcus2002,epicurus1994}. But their content is a claim about structure: that there is a determinate boundary between what is and is not within one's power, and that wisdom is the accurate tracking of that boundary, not a feat of self-overcoming. ``Taking the shape of the vessel'' read in the Stoic key is not a yielding of the self but a conforming to necessity --- bending where the structure of the situation genuinely admits no other course, and not bending where it does. This is exactly the de-cultivational reading the paper applies throughout: take a virtue that the tradition delivers as an inner discipline, and recover the structural claim inside it --- here, that the water-virtue's yielding is calibrated to the real joints of the situation, not an indiscriminate softness. Stoicism thus serves less as a source than as a worked example of how a seemingly contemplative doctrine is to be read for its structure, and we record it in that role.

% ---- Summary ----
\subsection{Summary: five challenges, and the home of the unanswered}
\label{p20:sub:review-summary}

The survey set the water-virtues against the secular ethical traditions and let the agreements and the conflicts stand without forcing a synthesis. What it yields is not a doctrine but a debt-sheet of five standing challenges, each routed to the place where the paper discharges it.

\emph{Challenge 1 (non-contention is irrational).} From game theory: non-contention is dominated except within a structure of repetition, reputation, and credible commitment. Routed to the survival cluster of the dialectical core (\xref{p20:sub:dialectic-survival}).

\emph{Challenge 1$'$ (non-contention is suicidal).} From Hobbes: absent an enforcing power, the non-contender is not exploited but eliminated. The graver form of the same demand, routed to the same place (\xref{p20:sub:dialectic-survival}), where it is answered not by erecting a sword but by the shaping of an ineliminable asymmetry.

\emph{Challenge 2 (taking-the-shape dissolves the subject).} From feminism and psychoanalysis: formless yielding may be the water-virtue or the self-erasure of a subject who has given way on its desire, and the figure of water does not distinguish them. Routed to the rigid-lower-bound cluster (\xref{p20:sub:dialectic-bound}).

\emph{Challenge 2$''$ (dwelling-below is slave morality).} From Nietzsche, anticipated in the political economy of the low place (\xref{p20:sec:lowliness}) and deferred: the descent the water-virtue praises may be the ressentiment-bred valorisation of weakness. Routed to the same cluster (\xref{p20:sub:dialectic-bound}), where it is answered by the decomposition that separates a generative descent from a slavish one.

\emph{Challenge 3 (formlessness is lawlessness).} From Kant --- and already answered, in place (\xref{p20:ssub:review-kant}), by reading the water-good as an imperfect duty: lawful in its ground, open in its form. It is collected here only for completeness; the dialectical core inherits the other four.

Two of the five thus belong to the cluster of survival (1, 1$'$) and two to the cluster of the rigid lower bound (2, 2$''$); the fifth is discharged. But the survey closes by naming something the challenges, and indeed the whole apparatus of criteria, structurally cannot reach. Every tradition canvassed argues about the \emph{skeleton} of the water-good --- the structural conditions under which a relation's circulation is generative rather than extractive, the criterion that sorts the good descent from the bad. None of them, and no structural criterion whatever, can reach the \emph{flesh}: the fact that this generative relation is \emph{this} one, with \emph{this} irreplaceable other, and not merely a well-formed instance of the generative type. The doctrines settle the type; the token --- the non-substitutable particularity in which the whole meaning of an intimate relation finally resides --- falls outside their reach by construction, and will be shown (\xref{p20:sub:formal-flesh}) to fall there for an exact formal reason. We carry forward, then, not only the four live challenges but this division of labour: the survey has mapped the skeleton, and marked the place where the flesh begins.
